\documentclass[11pt]{article}
\usepackage[a4paper,margin=27mm]{geometry}
\usepackage{amsmath,amssymb,amsthm,mathtools,xcolor,booktabs,array}
\newcommand{\PROVED}{\textcolor{green!40!black}{\textbf{PROVED}}}
\newcommand{\OPEN}{\textcolor{red!70!black}{\textbf{OPEN}}}
\newcommand{\AUDIT}{\textcolor{orange!70!black}{\textbf{AUDIT}}}
\newtheorem{theorem}{Theorem}
\newtheorem{lemma}{Lemma}
\newtheorem{corollary}{Corollary}
\title{The uniform all-depth compact-frequency tail for URGT$_1$ (v127)}
\date{14 August 2026}

\begin{document}
\maketitle

Put $L=\log N$ and
\begin{equation}
 B_{N,k}(\tau):=rac1{L^k}
 \sum_{m\le N}\frac{\Lambda_k(m)}{\sqrt m}\,m^{-i\tau}.
                                                                    \tag{1}
\end{equation}
The full ordered-pair polynomial is
$\overline{A_N(\tau)}B_{N,k}(\tau)$, and the leakage polynomial is
obtained by deleting its integral labels.  Since all coefficients are
nonnegative,
\begin{equation}
 |\Phi_{N,k}(\tau)|
 \le A_N(0)B_{N,k}(0).                                    \tag{2}
\end{equation}
Chebyshev's bound $\psi(x)\le Cx$ and partial summation give
\begin{equation}
 A_N(0)=\sum_{n\le N}\frac{\Lambda(n)}{\sqrt n}le C_0\sqrt N.
                                                                    \tag{3}
\end{equation}

\begin{lemma}[Logarithmic-shell convolution bound --- \PROVED]
There is an absolute constant $C_1$ such that, for $1\le k\le L$,
\begin{equation}
 \sum_{m\le N}\frac{\Lambda_k(m)}{\sqrt m}
 \le
 C_1^{k+1}\sqrt N
 \binom{\lfloor L\rfloor+k}{k-1}.                         \tag{4}
\end{equation}
Consequently
\begin{equation}
 B_{N,k}(0)
 \le\frac{C_1^{k+1}\sqrt N}{L}
 \left(\frac{2e}{k-1}\right)^{k-1}qquad(2\le k\le L),   \tag{5}
\end{equation}
with the $k=1$ case bounded separately by $C_0\sqrt N/L$.
\end{lemma}

\begin{proof}
On logarithmic shell $j\le\log n<j+1$, partial summation from
$\psi(x)\le Cx$ gives
\begin{equation}
 \sum_{e^j\le n<e^{j+1}}\frac{\Lambda(n)}{\sqrt n}
 \le C_1e^{j/2}.                                          \tag{6}
\end{equation}
Expand $\Lambda_k$ as the ordered $k$-fold Dirichlet convolution of
$\Lambda$.  A tuple $(n_1,\ldots,n_k)$ with product at most $N$ has shell
indices satisfying $j_1+\cdots+j_k\le L$.  For total shell index $q$, the
number of nonnegative shell compositions is
$\binom{q+k-1}{k-1}$.  Hence
\begin{align*}
 \sum_{m\le N}\frac{\Lambda_k(m)}{\sqrt m}
 &\le C_1^k
 \sum_{q\le L}e^{q/2}\binom{q+k-1}{k-1}\\
 &\le C_1^{k+1}e^{L/2}
       \binom{\lfloor L\rfloor+k}{k-1}.
\end{align*}
The last inequality follows by comparing successive terms backwards; the
factor $e^{-1/2}$ makes the resulting geometric majorant absolute.  This
is (4).  For $k\le L$,
\[
 \binom{\lfloor L\rfloor+k}{k-1}
 \le\left(\frac{e(L+k)}{k-1}\right)^{k-1}
 \le\left(\frac{2eL}{k-1}\right)^{k-1}.
\]
Division by $L^k$ proves (5).
\end{proof}

\begin{lemma}[Very large depths --- \PROVED]
There is an absolute $C_2$ such that
\begin{equation}
 B_{N,k}(0)\le N\left(\frac{C_2}{L}\right)^k              \tag{7}
\end{equation}
for every $k\ge1$.  Therefore
\begin{equation}
 \lim_{N\to\infty}
 \sum_{k>L}
 \left(\frac L N\right)^2
 |\Phi_{N,k}(\tau)|^2=0                                   \tag{8}
\end{equation}
uniformly in real $\tau$.
\end{lemma}

\begin{proof}
For every ordered tuple with $n_1\cdots n_k\le N$,
$1\le N/(n_1\cdots n_k)$.  Hence
\begin{align*}
 \sum_{m\le N}\frac{\Lambda_k(m)}{\sqrt m}
 &\le N
 \sum_{n_1,\ldots,n_k\ge2}
 \prod_{r=1}^k\frac{\Lambda(n_r)}{n_r^{3/2}}\\
 &=N\left(\sum_{n\ge2}\frac{\Lambda(n)}{n^{3/2}}\right)^k.
\end{align*}
This proves (7).  Combine (2)--(3) and (7):
\[
 \left(\frac L N\right)^2|\Phi_{N,k}(\tau)|^2
 \le C_0^2NL^2\left(\frac{C_2}{L}\right)^{2k}.
\]
For $k>L$, the geometric sum is
$O(NL^2(C_2/L)^{2L})=o(1)$, proving (8).
\end{proof}

\begin{theorem}[Uniform polynomial tail gate --- \PROVED]
For every fixed $R>0$,
\begin{equation}
 \boxed{
 \lim_{K\to\infty}\limsup_{N\to\infty}
 \sup_{|\tau|\le R}
 \left(\frac{\log N}{N}\right)^2
 \sum_{k>K}|\Phi_{N,k}(\tau)|^2=0.}                       \tag{9}
\end{equation}
Consequently the fixed-depth limits of v126 may be summed, and
\begin{equation}
 \left(\frac{\log N}{N}\right)^2W_N(\tau)
 \longrightarrow
 C_{\rm W}\frac1{(\tau^2+\frac14)^2},
 \qquad
 C_{\rm W}=\sum_{k\ge1}\frac1{((k-1)!)^2},               \tag{10}
\end{equation}
locally uniformly in $\tau$.
\end{theorem}

\begin{proof}
For $K<k\le L$, equations (2)--(3) and (5) give
\begin{equation}
 \left(\frac L N\right)^2|\Phi_{N,k}(\tau)|^2
 \le C_3^{2k}
 \left(\frac{2e}{k-1}\right)^{2k-2}.                     \tag{11}
\end{equation}
The series on the right is summable and its tail tends to zero independently
of $N$ and $\tau$.  Depths $k>L$ are $o(1)$ by Lemma 2.  This proves (9).
Fixed-depth convergence (5)--(6) of v126 and dominated convergence then
give (10).
\end{proof}

\begin{center}
\begin{tabular}{>{\raggedright\arraybackslash}p{.58\linewidth}
                >{\centering\arraybackslash}p{.14\linewidth}
                >{\raggedright\arraybackslash}p{.18\linewidth}}
\toprule
Statement & Status & Location \\
\midrule
Log-shell factorial bound & \PROVED & (4)--(6) \\
Superexponential very-large-depth tail & \PROVED & (7)--(8) \\
All-depth compact-frequency gate & \PROVED & (9) \\
Locally uniform all-depth profile & \PROVED & (10) \\
Uniform source-order compactness & \OPEN & v126, Lemma 2 \\
Finite-$N$ audit, URGT$_1$, $G$, row (d) & \OPEN & after compactness \\
\bottomrule
\end{tabular}
\end{center}

\end{document}
